\documentclass{article}

\usepackage{kantlipsum, xcolor, amsmath}
\usepackage{twoopt}
\usepackage{bm}

\title{\bf User Defined Commands}
\date{}

\renewcommand{\textheight}{650pt}

\begin{document}
	\maketitle
	
	\section{New Commands - Static} 
	\newcommand{\bibtex}{\textsc{Bib}\TeX~}
	\bibtex \\
	Typesetting \bibtex with the new command. \\
	
	\noindent
	\newcommand{\texnames}{Some logos: \TeX~\LaTeX}
	\texnames \\
	
	\newcommand{\creq}{
		{\color{blue}{$\mathsf{a^2=b^2+c^2}$}}	
	}
	Here is an equation \creq, more \creq
	
	\vspace{0.5cm}
	
	\section{New Commands - Dynamic}
	\newcommand{\sentence}[2]{% Max of 9 args.
		This is the \textit{#1} of sentence. 
		This is the #2 output.
	}
	\sentence{output}{second} \\
	\sentence{\textbf{output}}{second again}
	
	\newcommand{\kp}[1]{\kant[1-#1]}
	\kp{2}
	
	\newcommand{\texcmd}[1]{
		{
			\color{red!70}{\texttt{\detokenize{#1}}}		
		}
	}
	Trying the \verb|\texcmd| to produce \texcmd{\LaTeX}
	
%	\texcmd{\newcommand{\cmd}{def}}
	\noindent \\
	{\color{blue!50}{\verb|\newcommand{\cmd}{def}|}}
	
	\newcommand{\chain}[3]{
		% #1: y, #2: x, #3: u
		$
			\frac{d#1}{d#2}	= 
				\frac{d#1}{d#3} \cdot \frac{d#3}{d#2}
		$	
	}
	If $y=f(u)$ and $u=f(x)$, then \chain{y}{x}{u}
	
	\newcommand{\intx}[3]{
		$
			\int_{#3}^{\inf} #2\, d#1 	
		$
	}
	Integrating $x$: \intx{t}{x}{0} % {#1}{#2}{#3}
	
%	Trying to use $\intx{t}{x}{0}$

	\newcommand{\dd}[2]{
		\ensuremath{
			\frac{d#1}{d#2}		
		}	
	}
	Differentiating $x$ by time: $\dd{x}{t}$
	
	\subsection{Embedding Optional Arguments}
	
	\subsubsection{One as Optional}
	% \newcommand{\name}[1][def_first]{body}
	\newcommand{\bt}[1][Default]{
		{\color{blue}{#1}}
	}
	\bt \\
	\bt[New]
	
	\subsubsection{More than one, 1st being Optional}
	% \newcommand{\name}[args][def_first]{body}
	\newcommand{\txtfrac}[3][def]{
		\textit{#1/#2}	#3
	}
	\txtfrac{den}{three} \\
	\txtfrac[new]{den}{three again}
	
	\newcommand{\redtext}[2][]{
		\textcolor{red!#1}{#2}	
	}
	This is a \redtext{Red Text} Vs. New 
		\redtext[50]{Red Text}
		
	\subsection{Enhancing/Integrating User Commands}
	\newcommand{\x}{$x$\,}
	The x-command is now printing \x
	\newcommand{\xy}{
		\x $y$
	}
	Moreover, it's value can be integrated inside the 			new \xy
	
	\pagebreak
	
	\section{New Commands with Two Optional Arguments}
	\newcommandtwoopt{\twoopt}[3][A][B]{#1, #2, #3}
	\begin{description}
		\item[Mandatory with Default Optional]
			\twoopt{Third Argument}
		\item[Overriding First Optional Only]
			\twoopt[First New]{Third Argument}
		\item[Overriding both Optional]
			\twoopt[1st][2nd]{Still Three}
		\item[Empty First Argument]
			\twoopt[][Second]{Three Only}
		\item[Empty Second Argument]
			\twoopt[First][]{Third}
		\item[Empty both Optional]
			\twoopt[][]{Three Only}
	\end{description}
	
	\subsection{Renewing cmds with two optional args}
	\renewcommandtwoopt{\twoopt}[4][Q][P]{
		\underline{#1}
		\textcolor{#4}{#2}	
		\textit{#3}
	}
	\twoopt{Some input text}{green} \\[5pt]
	\twoopt[][\textsc{p}]{Some input text}{blue}\\[5pt]
	\twoopt[\texttt{Q}][]{Some input text}{red}
	
	\subsection{Omitting Defaults of Optional Arguments}
	\newcommandtwoopt{\txtforamt}[3][][]{
		#1 #2 #3	
	}
	\txtforamt{Some text} \\[5pt]
	\txtforamt[First Only]{Third Mandatory}
	
	\pagebreak
	
%	\renewcommand{\thesection}{\Alph{section}}
	\section{Redefining Commands}
%	\renewcommand{\thesubsection}{
%		\Alph{section}.\roman{subsection}
%	}
	\subsection{Altering the Built-in Commands}
	\renewcommand\tablename{Tabular}
	\renewcommand\thetable{\roman{table}}
	
	\begin{table}[h]
		\centering
		\begin{tabular}{|c|}
			\hline A \\ \hline 
		\end{tabular}
		\caption{Cell Table}
	\end{table}
	
	\noindent\\[5pt]
	\renewcommand{\textbf}[1]{
		\underline{\bfseries #1}
	}
	Some \textbf{New Text bf} surrounding text.
	
	\subsubsection{Be Careful Redefining Existing cmds}
	{
	The \textbackslash\texttt{vec} command was $\vec{v}$
	\renewcommand{\vec}[1]{ \mathbf{#1} }
	but now it's $\vec{v}$
	} \\
	So, we need to have a new command to exhibit
	\newcommand{\Vectorize}[1]{
		\bm{\vec{#1}}	
	}
	this $\Vectorize{v}$
	
	\subsection{Altering the User-defined Commands}
	\renewcommand{\xy}{xyz}
	The \verb|\xy| is now printing \xy 
	
%	\renewcommand{\thesection}{\arabic{section}}
	
	\pagebreak
	
	\section{Nesting Definitions of New Commands}
	
	\subsection{Using \textbackslash newcommand}
	\newcommand{\Father}{
		{% Group for Inner
		\newcommand{\Son}{
				\textcolor{blue}{Son Content}		
			}	
		Father contains \Son and \Son
		}% End group
	}
	\Father \\
	\Father \\
	\Father
	
	\begin{description}
		\item[Second Call] To call the outer command 					more than once you have to surround the 					inner commands with \{\}.
		\item[Scope of Inner] If the inner commands are 				surrounded by \{\}, then you cannot use them 			outside their scope, but if no brackets 					exist, then you can use them freely globally 			in the document.
	\end{description}
	
	\subsection{Using \textbackslash providecommand}
	\providecommand{\Father}{\textbf{Text}}
%	\newcommand{\Father}{\textit{Text}}
	\Father
	
	\noindent\\[5pt]
	\newcommand{\Outer}{
		\providecommand{\Inner}{
			\textcolor{orange}{Son Content}	
			\providecommand{\Inside}{
				\textcolor{red}{Grand Content}			
			}	
		}
		Outer contains \Inner and \Inside
	}
	\Outer \\
	\Outer
	
	\begin{description}
		\item[Provide Behavior] Allows for modular and 				conditional definitions of commands in \LaTeX. 				It's a feature often used in package development 		or when writing complex documents to ensure that 		commands are available when needed, but also to 			prevent errors due to command name clashes.
		\item[Scope of Inner] Note that the\verb|					\Inner| is accessible outside the parent 					container, \Inner and also \Inside
	\end{description}
	
	\subsection{Arguments for Nested Definitions}
	\newcommand{\father}[1]{
		\providecommand{\son}[1]{
			\textcolor{green!80}{##1}
			\providecommand{\grand}[1]{
				\textcolor{teal!80}{####1}	
				\providecommand{\fourth}[1]{
					\textcolor{orange!80}{########1}		
				}	
			}
		}
		\textcolor{red!80}{#1} contains 
			\son{Son Content} and \grand{Grand Content}
			and also \fourth{Fourth Level}!
	}
	\father{Father}
\end{document}